\documentclass[11pt,a4paper]{article}
\usepackage[utf8]{inputenc}
\usepackage[T1]{fontenc}
\usepackage[english]{babel}
\usepackage{lmodern}
\usepackage{amsmath,amssymb,amsthm}
\usepackage{booktabs}
\usepackage[margin=2.5cm]{geometry}
\usepackage[colorlinks=true,linkcolor=blue,citecolor=blue,urlcolor=blue]{hyperref}

\newtheorem{theorem}{Theorem}
\newtheorem{proposition}[theorem]{Proposition}
\newtheorem{lemma}[theorem]{Lemma}
\newtheorem{corollary}[theorem]{Corollary}
\newtheorem{conjecture}[theorem]{Conjecture}
\theoremstyle{remark}
\newtheorem{remark}{Remark}

\newcommand{\G}{\mathcal G}
\newcommand{\Ncal}{\mathcal N}
\newcommand{\dd}{\mathrm d}
\newcommand{\B}{\mathfrak B}
\setlength{\emergencystretch}{3em}

\title{The $d=N$ deficit and the throat ghost:\\
a common mechanism for the monopole instability of wormholes\\
in second-order scalar- and vector-tensor theories}
\author{Iván\\[2pt]
\small Independent researcher\thanks{Code and data: \href{https://github.com/iivanlopez/gusanos}{github.com/iivanlopez/gusanos}. Symbolic and numerical computations were carried out with the assistance of Claude (Anthropic).}}
\date{October 2026}

\begin{document}
\maketitle

\begin{abstract}
We study the monopole ($l=0$) sector of static, spherically symmetric wormhole throats with two asymptotically flat ends in the most general shift-symmetric second-order theories not already excluded by the Horndeski no-go theorems: quartic beyond-Horndeski (GLPV) theory, quintic Horndeski, the scalar linear in time $\phi=qt+\psi(\ell)$, and quartic generalized Proca. We show that the monopole obeys a common mechanism with two branches. (i)~If the throat is supported by a \emph{radial} field (a static scalar or the component $A_1$ of the vector), the trial function $\delta r=N$ (the areal radius rescaled by the lapse) satisfies all constraints and has energy exactly
$E[N]=-2\int N^3\G\,[(\ln r/N)']^2\,\dd\ell$, with no free functions, where $\G$ is the kinetic coefficient of the odd sector: if the latter is ghost-free, the monopole has a tachyon. (ii)~If it is supported by a \emph{temporal} field (the scalar $\phi=qt$ or the component $A_0$), the energy of the constraints can pay for the deficit, but the flaring-out condition $r''(0)>0$ fixes the sign of the kinetic energy at the throat: both reduced kinetic actions take the common form $T=\int[-(2\G/N)\dot d^2+k((\G\dot d/N)')^2]$ with $k(0)=-\kappa_0/(\G(0)r''(0))$, $\kappa_0>0$, so that there is a monopole ghost if $\G(0)>0$ and a tensor ghost if $\G(0)<0$ (for the temporal vector, $2\G$ is the coefficient $q_T$ of~\cite{DeFeliceProca}). For the scalar linear in time without a cubic term we further prove a stronger existence result: $G_4$ changes sign in every throat. We check that slow rotation, electric charge, a perfect fluid and an elastic solid do not change the sign of the deficit in the even sector, and that the exact Bakopoulos--Charmousis--Kanti throats have a strongly coupled monopole ($K\equiv0$) with $E[N]<0$. The only case not closed is the cubic galileon with a radial scalar flux ($G_3\neq0$, $\psi'\neq0$): we give two integral identities, the closed form of the principal symbol of the monopole, two new lemmas (the throat is a minimum of $NG_4$, and the tail must have a decreasing effective density) and a scan of about 3600 backgrounds in which none has a locally well-posed, ghost-free monopole.
\end{abstract}

\tableofcontents

%======================================================================
\section{Introduction}

A static traversable wormhole requires, in general relativity (GR), matter violating the null energy condition at the throat~\cite{MorrisThorne}. The canonical example, the Ellis--Bronnikov throat~\cite{Ellis,Bronnikov} supported by a phantom scalar, is unstable in the radial sector~\cite{GonzalezGuzmanSarbach}. Second-order scalar-tensor theories offer an alternative: the null-condition violation may be produced by a healthy scalar, or by a curvature coupling as in Einstein--scalar--Gauss--Bonnet~\cite{KKK,Canate}. In Horndeski theory, however, there are no static, spherically symmetric, asymptotically flat throats free of ghosts and gradient instabilities~\cite{Rubakov,EvseevMelichev}. The proof relies on a structure that ceases to exist beyond Horndeski~\cite{Franciolini}, and Mironov, Rubakov and Volkova (MRV) constructed throats in the quartic theory of Gleyzes--Langlois--Piazza--Vernizzi (GLPV)~\cite{GLPV} that satisfy the high-energy conditions in both parity sectors~\cite{MRV1,MRV2,MRV3}. The monopole sector has been treated as one more part of the even sector, and the complete analysis of Mironov and Volkova~\cite{MironovVolkova} assumes that it gives no additional conditions.

This work summarizes a programme (Phases~2f--2m of the repository~\cite{repo}) that studies precisely this sector, motivated by a hybrid object: a star with a throat at its centre. The result is negative and, above all, \emph{structural}: the monopole of every throat of these theories fails through one and the same mechanism, which depends only on the two-ended geometry and on the kinetic coefficient of the odd (or tensor) sector. Its form depends on which field supports the throat:
\begin{itemize}
\item \textbf{radial field} (a static scalar, or the radial component $A_1$ of a vector): a universal trial function, $\delta r=N$, has negative energy with no free functions (\S\ref{sec:radial});
\item \textbf{temporal field} (a scalar linear in time, or the component $A_0$): the energy of the constraints can compensate that deficit, but the flaring-out of the throat flips the sign of the kinetic energy (\S\ref{sec:temporal}).
\end{itemize}
The natural ways out (rotation, charge, matter, the quintic term, exact throats obtained by disformal transformation) do not change the conclusion (\S\ref{sec:exits}). One case remains without a theorem, the cubic galileon with radial flux, for which we give exact lemmas and numerical evidence (\S\ref{sec:cubic}).

\paragraph{Conventions.} Signature $(-,+,+,+)$, $X=-\tfrac12(\partial\phi)^2$, $G=c=1$ up to factors of $16\pi$ absorbed into the functions of the theory. The background is
\begin{equation}
\dd s^2=-N(\ell)^2\dd t^2+\dd\ell^2+r(\ell)^2\dd\Omega^2,\qquad N=e^{\Phi},
\label{eq:background}
\end{equation}
with $\ell\in\mathbb R$ the proper length, a throat at $\ell=0$ ($r'(0)=0$, $r''(0)>0$) and $r\to\infty$, $N\to1$ at both ends. The $l=0$ perturbations are $\delta g_{tt}=-2N^2a$, $\delta g_{t\ell}=\beta$, $\delta g_{\ell\ell}=2c$, $\delta r=d$.

%======================================================================
\section{Theories and structure of the monopole}\label{sec:framework}

\paragraph{Scalar.} Quartic GLPV theory~\cite{GLPV,Kobayashi},
\begin{equation}
\mathcal L=G_2-G_3\Box\phi+G_4R+G_{4X}\big[(\Box\phi)^2-\phi_{\mu\nu}^2\big]+F_4\,\epsilon^{\mu\nu\rho}{}_{\sigma}\epsilon^{\mu'\nu'\rho'\sigma}\phi_\mu\phi_{\mu'}\phi_{\nu\nu'}\phi_{\rho\rho'},
\label{eq:GLPV}
\end{equation}
with $G_i(\phi,X)$ and $F_4(\phi,X)$, and its quintic Horndeski extension with $G_5$. We define
\begin{equation}
\G=G_4-2XG_{4X}-4X^2F_4\qquad\big(\text{quintic: }\G=G_4-2XG_{4X}+X(\Phi'\phi'G_{5X}+G_{5\phi})\big).
\label{eq:Gcal}
\end{equation}
$\G$ is the kinetic coefficient of the odd sector~\cite{KMS1} (for $\G>0$ the odd sector has no ghosts) and also, as we shall see, that of tensor waves in the slicing $\phi=$~const.

\paragraph{Vector.} Quartic generalized Proca~\cite{Heisenberg},
\begin{equation}
\mathcal L=G_2(X,F)+G_3(X)\nabla_\mu A^\mu+G_4(X)R+G_{4X}\big[(\nabla_\mu A^\mu)^2-\nabla_\rho A_\sigma\nabla^\sigma A^\rho\big],
\end{equation}
with $A_\mu=(A_0,A_1,0,0)$, $X=-\tfrac12A_\mu A^\mu$ and $F=-\tfrac14F_{\mu\nu}F^{\mu\nu}$. Here $\G=G_4-2XG_{4X}$.

\paragraph{Machinery.} For each theory we derive symbolically the reduced spherical action and its quadratic $l=0$ expansion with general functions of the theory (for the even sector of Horndeski, see~\cite{KMS2}), and we check it with exact Noether identities in rational arithmetic, with exact gauge modes on on-shell jets, and with known solutions: Schwarzschild, Reissner--Nordström, linearized Kerr, Hartle--Thorne, the Lü--Pang black hole~\cite{LuPang} to 40 digits and the stealth black hole of Babichev--Charmousis~\cite{BabichevCharmousis} to 50 digits.

\paragraph{Structure of the monopole with a static scalar.} In unitary gauge $\delta\phi=0$ (Phases~2f and 2j), exactly,
\begin{equation}
L_2=-F\,W+L_{\rm rest}(c,d),\qquad F=N^2a'+\dot\beta,\qquad W=\gamma c+\alpha_1d'+\alpha_0d,\qquad \alpha_0=-\Phi'\alpha_1 .
\label{eq:structure}
\end{equation}
The lapse and shift enter only through the multiplier $F$, which imposes the momentum constraint $W=0$; $c$ is auxiliary; after eliminating it a single field $d$ remains, with $S_{\rm red}=\int\dd t\,\dd\ell\,[K\dot d^2-Pd'^2-Ud^2]$. The identity $\alpha_0=-\Phi'\alpha_1$ is exact in the quartic and in the quintic theory.

%======================================================================
\section{Radial field: the $d=N$ deficit}\label{sec:radial}

\begin{theorem}[monopole deficit; Phases 2g, 2j, 2m]\label{thm:dN}
Consider a background~\eqref{eq:background} with a static scalar $\phi(\ell)$ of quartic GLPV~\eqref{eq:GLPV} or quintic Horndeski theory, or with a purely radial vector ($A_0=0$) of quartic generalized Proca, with arbitrary functions of the theory. The configuration
\begin{equation}
d=N f,\qquad c=0,\qquad \delta\phi=0\ (\text{resp. }\delta A=0)
\end{equation}
satisfies all constraints, and its static energy is, pointwise and up to total derivatives,
\begin{equation}
E[Nf]=-2\int N^3\,\G\,\big(f'-\sigma f\big)^2\,\dd\ell,\qquad \sigma\equiv\Big(\ln\frac rN\Big)' .
\label{eq:EN}
\end{equation}
In particular, $E[N]=-2\int N^3\G\sigma^2\,\dd\ell<0$ if $\G>0$. If moreover the monopole is well posed ($K,P>0$), there is a mode with $\omega_0^2\le E[N]/\int KN^2\,\dd\ell<0$.
\end{theorem}

\begin{proof}[Sketch]
Since $\alpha_0=-\Phi'\alpha_1$, the constraint reads $W=\gamma c+\alpha_1N(d/N)'$, which vanishes for $c=0$ and $d\propto N$ (for constant $f$; for general $f$ it is corrected by $c$). The energy of the trial function is the part of $L_{\rm rest}$ containing only $d$, whose coefficients, on shell, depend only on $\G$; the identity~\eqref{eq:EN} is checked symbolically off shell. In the quintic theory, the reduction to the (unique) canonical form gives the Kobayashi--Motohashi--Suyama coefficient~\eqref{eq:Gcal}. For the radial vector, the background is that of Horndeski with $\phi'=A_1$, and in the gauge $\delta A_1=0$ the component $\delta A_0$ couples only to velocities. Finally, $\sigma\not\equiv0$ because $r/N\to\infty$ at both ends, and $d=N\to1$ is normalizable because $K$ is integrable.
\end{proof}

\begin{remark}
The Rayleigh quotient contains neither $G_2$, $G_3$ nor the normalization of the scalar: the theorem holds in strong and weak field, for any mass. In weak field ($|\Phi'|r\ll1$) the tachyon is \emph{fast}: $\omega_0^2\lesssim-O(1)/(Mr_0)$ (Phase~2f). For a throat with the mass of the Sun and $r_0\sim1$~km, the growth time is of milliseconds.
\end{remark}

\paragraph{Frame invariance.} With a static scalar, a disformal transformation $\bar g_{\mu\nu}=C(X)g_{\mu\nu}+D(X)\phi_\mu\phi_\nu$ leaves the trial function invariant ($\delta X=0$ and the lapse does not change), so its energy is the same in both frames (Phase~2j). This carries the theorem over to the exact Bakopoulos--Charmousis--Kanti (BCK) throats~\cite{BCK}, built by disformal transformation from the Lü--Pang black hole (\S\ref{sec:BCK}).

%======================================================================
\section{Temporal field: the throat ghost}\label{sec:temporal}

If the scalar depends linearly on time, $\phi=qt+\psi(\ell)$, the metric can be static in shift-symmetric theories~\cite{BabichevCharmousis}. Unitary gauge then fixes the \emph{time}, and the structure~\eqref{eq:structure} is lost: the lapse and shift of the slicing $\phi=$~const become algebraic auxiliaries with their own quadratic terms, and the trial function $d=N$ no longer satisfies the constraints. The same happens with the temporal component $A_0$ of a vector, which in Stückelberg language is a $\dot\pi$.

\subsection{The scalar linear in time}

With $\tau=t+\psi/q$ the scalar is $\phi=q\tau$, and in the ADM variables of the slicing (lapse $\Ncal$, shift $v$, proper length $\rho$) the Lagrangian is that of the effective theory~\cite{GLPV}
\begin{equation}
L=A_2+A_3K+A_4\big(K^2-K_{ij}K^{ij}\big)+B_4R^{(3)},\qquad A_4=-G_4+2XG_{4X}+4X^2F_4=-\G,\quad B_4=G_4 ,
\end{equation}
with $A_2=G_2$ and $A_3=yG_3-\int_0^yG_3\,\dd y$, $y=\sqrt{2X}$. The tensor speed in the slicing is $c_T^2=B_4/\G$. For stationary configurations we obtain in closed form (Phase~2k; the difference with the covariant Lagrangian has an identically vanishing Euler--Lagrange operator)
\begin{equation}
\mathcal L_{\rm st}=\Ncal r^2A_2-A_3(r^2v)'+\frac{2A_4}{\Ncal}\,r'(v^2r)'+\Ncal B_4\big(2-4rr''-2r'^2\big),
\end{equation}
with the momentum constraint $rA_3'=2v\,(Qr')'$, $Q\equiv2A_4/\Ncal$. Without a cubic term ($A_3'\equiv0$ on the background), a throat forces $v\equiv0$ ($\psi'=0$) or $A_4\equiv0$.

\begin{theorem}[existence no-go; Phase 2k]\label{thm:beta}
Consider a static throat with $\phi=qt+\psi$ in shift-symmetric quartic GLPV theory, with $v\equiv0$ or $A_3'\equiv0$, two asymptotically flat ends ($r'/r^2\to0$ and $\beta'/(r\beta)\to0$) and $G_4\to G_{4\infty}\neq0$. Then $\beta\equiv\Ncal G_4$ changes sign; since $\Ncal>0$, $G_4$ changes sign on the background.
\end{theorem}

\begin{proof}
Eliminating $A_2$ between the first integral and the radial equation, $\beta$ satisfies, with no free function,
\begin{equation}
\Big(\frac{\beta'}{r}\Big)'+\frac{K\beta}{r^3}=0,\qquad K=1+rr''-r'^2 .
\label{eq:beta}
\end{equation}
If $\beta$ does not vanish, dividing by $\beta$, integrating by parts and using $K/r^3=(r'/r^2)'+(1+r'^2)/r^3$,
\begin{equation}
0\le\int_{\mathbb R}\frac{\beta'^2}{r\beta^2}\,\dd\rho=\Big[\frac{\beta'}{r\beta}\Big]-\int\frac{K}{r^3}\,\dd\rho=-\Big[\frac{r'}{r^2}\Big]-\int\frac{1+r'^2}{r^3}\,\dd\rho<0. \qedhere
\end{equation}
\end{proof}

Neither reflection symmetry nor uniqueness of the throat is needed; what matters is the two-ended topology (with a regular centre, $r(0)=0$ and the integral does not exist). Where $B_4=G_4$ vanishes, $c_T^2=0$ (strong coupling); on one side of the zero there is a tensor gradient instability. In GR ($\beta=N$), \eqref{eq:beta} is the isotropic-pressure condition.

\paragraph{The monopole with $\phi=qt$.} The time-dependent quadratic action (exact up to total derivatives) contains the lapse only in the static part and the shift only in the kinetic part. Eliminating them,
\begin{equation}
T[\dot d]=\int\Big[Q\dot d^2+\frac{r\,[(Q\dot d)']^2}{(Qr')'}\Big]\dd\rho,\qquad
V[d]=\int\Big[\frac{S^2}{2\kappa}-2\beta r^2\Big[\Big(\frac dr\Big)'\Big]^2\Big]\dd\rho,
\label{eq:Tqt}
\end{equation}
with no gyroscopic terms. The family $\delta X=0$, the analogue of the trial function $d=N$, has energy $-2\int\beta r^2[(d/r)']^2$: the same form as~\eqref{eq:EN}, with $\G\to G_4$, and Theorem~\ref{thm:beta} forbids precisely $\beta>0$.

\subsection{The temporal vector}

With $A_1=0$ and $G_3=0$ (Phase~2m), the exact structure is $L_2=\kappa F^2-FW+L_{\rm rest}$ with $\kappa=U^2r^2(G_{2F}+2FG_{2FF})/(2N^5)$ ($U$ is $A_0$ on the background). The trial function $d=N$ does not satisfy the constraint, and eliminating $F$ and the algebraic perturbation of $A_0$ the energy of the perturbed electric field appears,
\begin{equation}
E[N]=\int\Big[-2N^3G_4\sigma^2+\frac{2N}{g}\big(2G_{4X}N(U/N)'-G_{2F}U'\big)^2+\frac{J_N^2}{4\tilde m}\Big]\dd\ell,\qquad g=G_{2F}+2FG_{2FF},
\end{equation}
which in the six regular throats we constructed exceeds the deficit by a factor of 1.2 to 4.6. But the throat imposes an exact local condition,
\begin{equation}
r''(0)=-\frac{X\Theta}{2r_0\G},\qquad \tilde c(0)=-\frac{N_0\Theta}{2}=\frac{N_0r_0\G}{X}\,r''(0),\qquad \Theta\equiv G_{2X}r_0^2+2G_{4X},
\end{equation}
where $\tilde c$ is the coefficient of the auxiliary radial component $\tilde w=\delta A_1+U\beta/N^2$. The complete reduction in the gauge $c=0$ gives
\begin{equation}
S_{\rm red}=\int\Big[-\frac{2\G}{N}\dot d^2+k_2\Big(\Big(\frac{\G\dot d}{N}\Big)'\Big)^2-V(d)\Big],\qquad k_2=-\frac{4r^2N^4}{U^2\tilde c}.
\end{equation}

\subsection{The common sign rule}

\begin{theorem}[throat ghost]\label{thm:ghost}
In the two temporal cases without flux ($\phi=qt$ with $\psi'=0$; vector with $A_1=0$, $G_3=0$) the kinetic energy of the monopole is
\begin{equation}
T=\int\Big[-\frac{2\G}{N}\,\dot d^2+k\,\Big(\Big(\frac{\G\dot d}{N}\Big)'\Big)^2\Big]\dd\ell,\qquad k(0)=-\frac{\kappa_0}{\G(0)\,r''(0)},\quad \kappa_0>0,
\label{eq:Tcommon}
\end{equation}
with $k=4r/(Qr')'$, $\kappa_0=2N_0r_0$ for the scalar and $k=k_2$, $\kappa_0=4N_0^3Xr_0/U_0^2$ for the vector. Hence, near a throat ($r''(0)>0$) both kinetic terms have the sign of $-\G(0)$: if the tensor waves of the slicing are not ghosts ($\G(0)>0$), every $l=0$ perturbation supported near the throat has negative kinetic energy; if $\G(0)<0$, the ghost is in the tensor sector.
\end{theorem}

\begin{proof}
For the scalar, $Q=2A_4/N=-2\G/N$ turns~\eqref{eq:Tqt} into~\eqref{eq:Tcommon} with $k=4r/(Qr')'$, and at the throat $(Qr')'=Qr''$. For the vector, it suffices to substitute $\tilde c(0)$ into $k_2$. Both substitutions are checked symbolically (\texttt{fase2n/regla\_comun.py}).
\end{proof}

With $G_4=1$, $(Qr')'\propto\rho+p$ is the null-condition combination, and~\eqref{eq:Tcommon} is the kinetic energy of a fluid: the theorem is the general version of ``phantom fluid $\Rightarrow$ ghost''. The refinement given by the principal symbol at the throat (Phase~2k \S12) is that, with $\G(0)>0$, either the radial speeds are imaginary (ill-posed problem) or the whole sector has the opposite sign to the gravitons. With untilted causal cones at the throat, the slicing criterion of Babichev et al.~\cite{BCEL1,BCEL2} does not avoid it.

%======================================================================
\section{Ways out that do not work}\label{sec:exits}

\subsection{Charge and slow rotation}

A radial electric field (``charge without charge'') adds no degrees of freedom at $l=0$ and contributes a favourable pressure, $E[N]=\int[-2N^3\G\sigma^2+8q_e^2N^3/r^4]$. But in a symmetric throat the coefficient of the equation for $c$ no longer vanishes, $e_0=4q_e^2N_0/r_0^3$, which forces $\delta r(0)=0$ and produces a fall to the centre, $U\simeq-C/\ell^2$ with $C/P(0)=7.5$--$19$ (Phase~2g).

In slow rotation (Phase~2l), the Kaluza--Klein reduction of quartic GLPV gives exactly $L_\omega=\tfrac12\mathcal J\omega'^2$ with $\mathcal J=\tfrac23C^2\G/\sqrt{AB+b^2}$, the same $\G$. At fixed angular momentum, the $l=0$ sector at $O(\Omega^2)$ is a Routh problem with $L_{\rm eff}=L_{l=0}-\pi^2/(2\mathcal J)$, and
\begin{equation}
E_J[Nf]=\int\Big[-2N^3\G(f'-\sigma f)^2+\frac{9\pi^2N^3}{\G r^6}f^2\Big]\dd\ell .
\end{equation}
With $v_0$ the speed of the zero-angular-momentum observer at the throat, $E_J[N]/|E_0|=-1+\kappa v_0^2$: for Ellis--Bronnikov $\kappa=24/\pi^2$ and the sign only changes at $v_0=\pi/\sqrt{24}\approx0.64$; in 27 reconstructed throats, $v_c=0.59$--$1.74$. In the controlled regime ($v_0\le0.3$) rotation reduces $|E[N]|$ by at most 26\,\%. Moreover it reintroduces the charge obstruction at the throat, except in the subclass $(\G-\phi'^2\G_X)(0)=0$.

\subsection{Matter: fluid and solid}

With a perfect fluid (Phases~2g and 2h), the trial function with the fluid at rest gives $E=-2\int N^3[\G\sigma^2-16\pi\Gamma_1p]$, and in the Sun $16\pi\Gamma_1pr_0^2\sim10^{-20}$--$10^{-8}$: the solar hybrid cannot be saved. The dense hybrids that satisfied the necessary conditions with the stellar surface held fixed fail with a free surface: the \emph{compression-free} displacement $y$ has
\begin{equation}
E_{Q0}=-2\int_0^{\ell_s}N^3\G(y'-\sigma y)^2\,\dd\ell+E_{\rm out}\,y(\ell_s)^2+E_s\,y'(\ell_s)^2 ,
\end{equation}
with the first two terms negative. This is a theorem for soft surfaces ($P\propto(Y-Y_s)^m$, $1<m<2$); with a density jump, the best margin in about 3200 optimized backgrounds is $-0.22$. The speed of the fast fluid mode moreover diverges at the throat.

An elastic solid (Phase~2i; a field formulation equivalent to that of~\cite{KarloviniSamuelsson}) stabilizes the odd sector of the monopole (from $-0.34$ to $+0.06$ with $c_T^2=0.2$), but not the even one: in every region where the background satisfies the GR relation $\Phi'r'-r''=r(\tilde\rho+\tilde p_r)/(4\G)$, the uniform translation $d=r'C$, $\xi=-C$ does not strain the material and satisfies the constraint, with energy
$E_m=C^2N[\tilde p_t(r'^2+rr'')+\tilde\rho\,r^2N''/(2N)]$ \emph{without any elastic modulus}. The best margin in about 2900 backgrounds is $-0.209$.

\subsection{BCK throats}\label{sec:BCK}

In the frame of the seed (the exterior of the Lü--Pang black hole, doubled), $E[N]<0$ for all nine parameter sets tested ($-0.17$ to $-0.53$ in units of $M$), with $\G\ge1$. Furthermore, the kinetic coefficient of the monopole vanishes identically, $K\equiv0$, on that solution (not on generic quintic backgrounds), and $P<0$ in some region: the BCK monopole is strongly coupled and linear analysis does not apply.

%======================================================================
\section{The open case: cubic galileon with flux}\label{sec:cubic}

With $G_3\neq0$ and $\psi'\neq0$ the momentum constraint allows $v\neq0$ with $A_4\neq0$, and~\eqref{eq:beta} acquires a source:
\begin{equation}
\Big(\frac{\beta'}r\Big)'+\frac{K\beta}{r^3}=-\frac{(Qr^4u')'}{4r^3},\qquad u=\frac{v^2}{r^2}.
\label{eq:betasource}
\end{equation}
A symmetric configuration has odd $v$. The health conditions of the background are $G_4>0$, $A_4<0$ (that is, $\G>0$, $Q<0$), $\Ncal>|v|$ and $G_4\to1$, $A_4\to-1$ at both ends.

\paragraph{There is no algebraic no-go.} Two exact identities of~\eqref{eq:betasource}:
\begin{align}
&\text{(I1)}\quad \int\frac{K\beta}{r^3}\,\dd\rho=\frac34\int\Big(\frac vr\Big)^2(Qr')'\,\dd\rho,\\
&\text{(I2)}\quad \int\frac1r\Big(\frac{\beta'}\beta+\frac{Qr^2u'}{8\beta}\Big)^2+\int\frac{1+r'^2}{r^3}=-\frac34\int\frac{Qu'r'}\beta+\frac1{64}\int\frac{Q^2r^3u'^2}{\beta^2}\qquad(\beta>0).
\end{align}
The last term of (I2) is positive: a sufficiently strong flux pays for the two-ended geometry, and there are explicit backgrounds with $G_4>0$ and $A_4<0$ along the whole axis.

\paragraph{Principal symbol.} The time-dependent ADM Lagrangian with shift is exact relative to the covariant one, and the principal symbol of the monopole has a closed form. At any point, the kinetic coefficient in the slicing $\tau$ is
\begin{equation}
s_2=-\frac{2\B}{r\,v^2r''^2},\qquad \B=-2\beta'(rK)'-Q\,w\,r''(r^4w')'+Q'\,r^4w'(r'w'-r''w),\qquad w=\frac vr ,
\label{eq:B}
\end{equation}
independent of the lapse. At the throat it depends only on Taylor data of the background:
\begin{equation}
\hat\Lambda=-\frac{\Ncal_0}{4R_0^3}\Big[Q_0^2R_0^4V_1^4+2Q_0R_0^3R_2V_1^2\beta_0-4\beta_0^2(1+R_0R_2)^2\Big],
\end{equation}
with $R_k=r^{(k)}(0)$, $V_k=v^{(k)}(0)$; the gradient is healthy if $\hat\Lambda/A_4>0$.

\begin{lemma}[the throat is a minimum of $\beta$; Phase 2n]\label{lem:minimum}
If $\beta_0>0$, $Q_0<0$ and the gradient of the monopole is healthy at the throat ($\hat\Lambda<0$), then $\beta''(0)>R_2\beta_0/(2R_0)>0$.
\end{lemma}
\begin{proof}
From~\eqref{eq:betasource} at $\rho=0$, $\beta''(0)=-(\beta_0/R_0^2)\,[1+R_0R_2+x/2]$ with $x=Q_0R_0^2V_1^2/\beta_0<0$. The bracket of $\hat\Lambda$ is $\beta_0^2[x^2+2R_0R_2x-4(1+R_0R_2)^2]$, positive with $x<0$ only if $-x>R_0R_2+\sqrt{R_0^2R_2^2+4(1+R_0R_2)^2}>R_0R_2+2(1+R_0R_2)$. Substituting, $\beta''(0)R_0^2/\beta_0>R_0R_2/2$.
\end{proof}

\begin{lemma}[the tail; Phase 2n]\label{lem:tail}
(a) For any $r(\rho)$, $r^3\big[(r''/r)'+Kr'/r^3\big]=(rK)'$; that is, $r'$ solves~\eqref{eq:beta} if and only if $rK$ is constant. (b) With the Misner--Sharp mass of the slicing, $r'^2=1-2m(r)/r$, one has $K=3m/r-m_r$ and
\begin{equation}
(rK)'=-4\pi\,r'\,r^3\,\frac{\dd\varepsilon}{\dd r},\qquad \varepsilon\equiv\frac{m_r}{4\pi r^2}.
\end{equation}
(c) Where the flux fades ($v,v',v''\to0$), $\B\to-2\beta'(rK)'$, so a healthy kinetic coefficient requires $\beta'\,\dd\varepsilon/\dd r<0$. If moreover the flux is negligible further out (where $\beta'/r=\int_\rho^\infty K\beta/r^3>0$), the effective density $\varepsilon$ must be decreasing.
\end{lemma}

Lemma~\ref{lem:tail} excludes, for any $A_4$ and lapse, all Ellis-type tails ($\varepsilon<0$ increasing towards zero, $(rK)'=-2r'/r^2$). Note that the kinetic condition of the flux-free regime, $(Qr')'>0$, can be satisfied in an Ellis tail with a tail profile of $A_4$; it is the high-wavenumber regime that excludes it. On the other hand, the affine dependence of $\B$ on $Q'$ prevents closing the problem with a linear identity $\int\mu\B$ with fixed weight: after integrating $Q'$ by parts, the integrand $Q\,C$ contains the term $-2\mu r^4r'w'w''$, of indefinite sign.

\paragraph{Numerical evidence.} We reconstruct backgrounds by choosing $(r,v,Q,\Ncal)$ and solving~\eqref{eq:betasource} for the bounded even solution, and we evaluate the principal symbol at 420--750 points per background:
\begin{center}
\small\begin{tabular}{p{6.2cm}ccc}
\toprule
family & backgrounds & $\beta>0$ & locally healthy \\
\midrule
Ellis, one or two fluxes, peaked $A_4$ (2k) & 872 & 56 & 0 \\
positive-energy tail, $\varepsilon>0$ on the whole axis (2k) & 400 (+ optim.) & 0 & 0 \\
Ellis core with positive-energy tail (2n) & 600 (+ 2160 optim.) & 34 & 0 \\
\bottomrule
\end{tabular}
\end{center}
The backgrounds fail either through a non-hyperbolic layer next to the throat (where $\B$ changes sign) or through a tail where the kinetic coefficient is negative. The class with an Ellis core and a tail of decreasing density, the one left by Lemma~\ref{lem:tail}, admits $\beta>0$ with a healthy tail (19 of 600 backgrounds), contrary to what the families with $\varepsilon>0$ at the throat suggested. Differential evolution in this class finds a background that is hyperbolic along the whole axis, with $\beta>0$ and healthy throat and tail, but with a ghost layer at $0.56\lesssim\rho/r_0\lesssim1.7$: there $s_2<0$ and $s_0>0$ simultaneously. The two characteristic speeds have opposite signs (their product is $s_0/s_2<0$), so the slicing $\tau$ is spacelike for the monopole and the negative energy is a ghost also according to the criterion of~\cite{BCEL2}. This is the region where the peak of $|A_4|$ supporting the throat unwinds ($Q'>0$) and the term affine in $Q'$ in~\eqref{eq:B} has the wrong sign.

\begin{conjecture}\label{conj:cubic}
Every static, symmetric, asymptotically flat throat of the cubic galileon with flux in shift-symmetric quartic GLPV theory has an ill-posed monopole or one with a ghost.
\end{conjecture}

%======================================================================
\section{Discussion}\label{sec:discussion}

\paragraph{The mechanism.} In both cases the cause is the two-ended topology. With a radial field, the throat pays with $\sigma=(\ln r/N)'\neq0$: the ratio $r/N$ has to grow towards both sides, and the trial function $d=N$ turns this into negative energy weighted by the same $\G$ that makes the odd sector healthy. With a temporal field, the integrated Sturm--Liouville potential of~\eqref{eq:beta}, $\int(1+r'^2)/r^3$, is positive, and the flaring-out $r''(0)>0$ ties the sign of the kinetic energy to that of $-\G(0)$. No term of the theory can compensate the former (it contains no free functions); the latter can only be compensated by flipping the tensor sector.

\paragraph{Relation to the literature.} Theorem~\ref{thm:dN} completes the conditions of MRV and of Mironov--Volkova~\cite{MRV3,MironovVolkova} in the monopole sector: their throats with healthy high-energy behaviour have a tachyon at $l=0$. For the scalar linear in time, the instability of hairy black holes~\cite{OgawaKobayashiSuyama} and the strong coupling of stealth solutions~\cite{deRhamZhang} have their analogue here in the sign change of $G_4$ (Theorem~\ref{thm:beta}). For generalized Proca, the result is consistent with the known instabilities of solutions with vector hair~\cite{HKMT,KaseOdd,DeFeliceCubic,GarciaSaenz} and with the observation that minimal vector wormholes require a ghost~\cite{Ganiyeva}.

\paragraph{What is not covered.} (i) The cubic galileon with flux (Conjecture~\ref{conj:cubic}) and the mixed vector case $A_0,A_1\neq0$, which is the same gap and in addition breaks time parity (there are gyroscopic terms). (ii) The quintic ADM reduction with $\phi=qt+\psi$. (iii) General degenerate DHOST theories outside the classes connected by disformal transformations with a static scalar. (iv) Relativistic rotation ($v_0\gtrsim0.6$), outside perturbation theory. (v) Matter coupled directly to the scalar, which could ``feel'' it throughout its volume and prevent the near-gauge translation.

\paragraph{Consequence for the star--throat hybrid.} The solar throat is extremely weak-field ($M/r_0\sim10^{-6}$): by Theorem~\ref{thm:dN} it has a tachyon with growth time $\sim\sqrt{Mr_0}$, of milliseconds, and the matter of the star does not stop it.

%======================================================================
\section{Conclusions}

\begin{enumerate}
\item With a radial field, the trial function $d=N$ gives $E[N]=-2\int N^3\G\sigma^2<0$ with no free functions in quartic GLPV, in quintic Horndeski and in radial generalized Proca: a tachyon if the odd sector is healthy.
\item With a temporal field without flux, the kinetic energy of the monopole takes the common form~\eqref{eq:Tcommon} and the throat fixes its sign: a monopole ghost if $\G(0)>0$, a tensor ghost if $\G(0)<0$. For the scalar, moreover, $G_4$ changes sign.
\item Slow rotation, charge, fluid and solid do not change the sign of the even sector; the BCK throats are strongly coupled.
\item The cubic galileon with flux remains open, with exact lemmas restricting the throat and the tail, and about 3600 backgrounds without any locally healthy monopole; the best one is hyperbolic, but has a ghost layer between the throat and the maximum of the flux.
\end{enumerate}

\appendix
\section{Reproducibility}
All computations are in the repository~\cite{repo}. The symbolic machinery and the checks of each phase in \texttt{fase2f/}--\texttt{fase2m/}, and the lemmas of \S\ref{sec:cubic} and the common sign rule in \texttt{fase2n/} (\texttt{lemas.py}, \texttt{regla\_comun.py}, \texttt{criba\_mixta.py}, \texttt{optimiza\_mixta.py}). The arXiv references were verified with its API (\texttt{fase2n/arxiv\_bib.py}). The reports of each phase are written in Spanish.

\begin{thebibliography}{99}

\bibitem{MorrisThorne} M.~S. Morris and K.~S. Thorne, \emph{Wormholes in spacetime and their use for interstellar travel: A tool for teaching general relativity}, Am. J. Phys. \textbf{56} (1988) 395.

\bibitem{Ellis} H.~G. Ellis, \emph{Ether flow through a drainhole: A particle model in general relativity}, J. Math. Phys. \textbf{14} (1973) 104.

\bibitem{Bronnikov} K.~A. Bronnikov, \emph{Scalar-tensor theory and scalar charge}, Acta Phys. Pol. B \textbf{4} (1973) 251.

\bibitem{GonzalezGuzmanSarbach} J.~A. González, F.~S. Guzmán and O. Sarbach, \emph{Instability of wormholes supported by a ghost scalar field. I. Linear stability analysis}, Class. Quant. Grav. \textbf{26} (2009) 015010, \href{https://arxiv.org/abs/0806.0608}{arXiv:0806.0608}.

\bibitem{Rubakov} V.~A. Rubakov, \emph{Can Galileons support Lorentzian wormholes?}, \href{https://arxiv.org/abs/1509.08808}{arXiv:1509.08808}.

\bibitem{EvseevMelichev} O.~A. Evseev and O.~I. Melichev, \emph{No static spherically symmetric wormholes in Horndeski theory}, Phys. Rev. D \textbf{97} (2018) 124040, \href{https://arxiv.org/abs/1711.04152}{arXiv:1711.04152}.

\bibitem{Franciolini} G. Franciolini, L. Hui, R. Penco, L. Santoni and E. Trincherini, \emph{Stable wormholes in scalar-tensor theories}, JHEP \textbf{01} (2019) 221, \href{https://arxiv.org/abs/1811.05481}{arXiv:1811.05481}.

\bibitem{GLPV} J. Gleyzes, D. Langlois, F. Piazza and F. Vernizzi, \emph{Healthy theories beyond Horndeski}, Phys. Rev. Lett. \textbf{114} (2015) 211101, \href{https://arxiv.org/abs/1404.6495}{arXiv:1404.6495}.

\bibitem{MRV1} S. Mironov, V. Rubakov and V. Volkova, \emph{Towards wormhole beyond Horndeski}, EPJ Web Conf. \textbf{191} (2018) 07014, \href{https://arxiv.org/abs/1811.05832}{arXiv:1811.05832}.

\bibitem{MRV2} S. Mironov, V. Rubakov and V. Volkova, \emph{More about stable wormholes in beyond Horndeski theory}, \href{https://arxiv.org/abs/1812.07022}{arXiv:1812.07022}.

\bibitem{MRV3} S. Mironov, V. Rubakov and V. Volkova, \emph{In hot pursuit of a stable wormhole in beyond Horndeski theory}, \href{https://arxiv.org/abs/2212.05969}{arXiv:2212.05969}.

\bibitem{MironovVolkova} S. Mironov and V. Volkova, \emph{Complete stability for spherically symmetric backgrounds in beyond Horndeski theory}, \href{https://arxiv.org/abs/2404.06297}{arXiv:2404.06297}.

\bibitem{repo} \emph{gusanos: star--wormhole hybrid} (2026), \url{https://github.com/iivanlopez/gusanos}.

\bibitem{Kobayashi} T. Kobayashi, \emph{Horndeski theory and beyond: a review}, Rept. Prog. Phys. \textbf{82} (2019) 086901, \href{https://arxiv.org/abs/1901.07183}{arXiv:1901.07183}.

\bibitem{KMS1} T. Kobayashi, H. Motohashi and T. Suyama, \emph{Black hole perturbation in the most general scalar-tensor theory with second-order field equations I: The odd-parity sector}, Phys. Rev. D \textbf{85} (2012) 084025, \href{https://arxiv.org/abs/1202.4893}{arXiv:1202.4893}.

\bibitem{KMS2} T. Kobayashi, H. Motohashi and T. Suyama, \emph{Black hole perturbation in the most general scalar-tensor theory with second-order field equations II: the even-parity sector}, Phys. Rev. D \textbf{89} (2014) 084042, \href{https://arxiv.org/abs/1402.6740}{arXiv:1402.6740}.

\bibitem{Heisenberg} L. Heisenberg, \emph{Generalization of the Proca Action}, JCAP \textbf{05} (2014) 015, \href{https://arxiv.org/abs/1402.7026}{arXiv:1402.7026}.

\bibitem{LuPang} H. Lü and Y. Pang, \emph{Horndeski gravity as $D\to4$ limit of Gauss--Bonnet}, Phys. Lett. B \textbf{809} (2020) 135717, \href{https://arxiv.org/abs/2003.11552}{arXiv:2003.11552}.

\bibitem{BabichevCharmousis} E. Babichev and C. Charmousis, \emph{Dressing a black hole with a time-dependent Galileon}, JHEP \textbf{08} (2014) 106, \href{https://arxiv.org/abs/1312.3204}{arXiv:1312.3204}.

\bibitem{BCK} A. Bakopoulos, C. Charmousis and P. Kanti, \emph{Traversable wormholes in beyond Horndeski theories}, JCAP \textbf{05} (2022) 022, \href{https://arxiv.org/abs/2111.09857}{arXiv:2111.09857}.

\bibitem{BCEL1} E. Babichev, C. Charmousis, G. Esposito-Farèse and A. Lehébel, \emph{Stability of a black hole and the speed of gravity waves within self-tuning cosmological models}, Phys. Rev. Lett. \textbf{120} (2018) 241101, \href{https://arxiv.org/abs/1712.04398}{arXiv:1712.04398}.

\bibitem{BCEL2} E. Babichev, C. Charmousis, G. Esposito-Farèse and A. Lehébel, \emph{Hamiltonian vs stability and application to Horndeski theory}, Phys. Rev. D \textbf{98} (2018) 104050, \href{https://arxiv.org/abs/1803.11444}{arXiv:1803.11444}.

\bibitem{OgawaKobayashiSuyama} H. Ogawa, T. Kobayashi and T. Suyama, \emph{Instability of hairy black holes in shift-symmetric Horndeski theories}, Phys. Rev. D \textbf{93} (2016) 064078, \href{https://arxiv.org/abs/1510.07400}{arXiv:1510.07400}.

\bibitem{deRhamZhang} C. de Rham and J. Zhang, \emph{Perturbations of stealth black holes in DHOST theories}, Phys. Rev. D \textbf{100} (2019) 124023, \href{https://arxiv.org/abs/1907.00699}{arXiv:1907.00699}.

\bibitem{HKMT} L. Heisenberg, R. Kase, M. Minamitsuji and S. Tsujikawa, \emph{Hairy black-hole solutions in generalized Proca theories}, Phys. Rev. D \textbf{96} (2017) 084049, \href{https://arxiv.org/abs/1705.09662}{arXiv:1705.09662}.

\bibitem{KaseOdd} R. Kase, M. Minamitsuji, S. Tsujikawa and Y.-l. Zhang, \emph{Black hole perturbations in vector-tensor theories: The odd-mode analysis}, JCAP \textbf{02} (2018) 048, \href{https://arxiv.org/abs/1801.01787}{arXiv:1801.01787}.

\bibitem{DeFeliceProca} A. De Felice, L. Heisenberg, R. Kase, S. Mukohyama, S. Tsujikawa and Y.-l. Zhang, \emph{Cosmology in generalized Proca theories}, JCAP \textbf{06} (2016) 048, \href{https://arxiv.org/abs/1603.05806}{arXiv:1603.05806}.

\bibitem{DeFeliceCubic} A. De Felice, R. Kase and S. Tsujikawa, \emph{Scrutinizing black hole stability in cubic vector Galileon theories}, \href{https://arxiv.org/abs/2409.15606}{arXiv:2409.15606}.

\bibitem{GarciaSaenz} S. Garcia-Saenz, A. Held and J. Zhang, \emph{Destabilization of black holes and stars by generalized Proca fields}, Phys. Rev. Lett. \textbf{127} (2021) 131104, \href{https://arxiv.org/abs/2104.08049}{arXiv:2104.08049}.

\bibitem{Ganiyeva} N. Ganiyeva, B.~J. Barros, Á. de la Cruz-Dombriz and F.~S.~N. Lobo, \emph{Vector wormholes as conduits for matter interaction}, Eur. Phys. J. C \textbf{85} (2025) 914, \href{https://arxiv.org/abs/2412.19408}{arXiv:2412.19408}.

\bibitem{KarloviniSamuelsson} M. Karlovini and L. Samuelsson, \emph{Elastic stars in general relativity: I. Foundations and equilibrium models}, Class. Quant. Grav. \textbf{20} (2003) 3613, \href{https://arxiv.org/abs/gr-qc/0211026}{arXiv:gr-qc/0211026}.

\bibitem{KKK} P. Kanti, B. Kleihaus and J. Kunz, \emph{Wormholes in dilatonic Einstein--Gauss--Bonnet theory}, Phys. Rev. Lett. \textbf{107} (2011) 271101, \href{https://arxiv.org/abs/1108.3003}{arXiv:1108.3003}.

\bibitem{Canate} P. Cañate, J. Sultana and D. Kazanas, \emph{Ellis wormhole without a phantom scalar field}, Phys. Rev. D \textbf{100} (2019) 064007, \href{https://arxiv.org/abs/1907.09463}{arXiv:1907.09463}.

\end{thebibliography}
\end{document}
